\documentclass[10pt,a4paper]{article}
\usepackage{amsmath,amssymb,amsthm,graphicx}
\usepackage{hyperref}
\usepackage{fontspec}
\title{Randomized average-case easiness yields subexponential simulations of promise polynomial hierarchy levels}
\author{Halley Goldberg, Valentine Kabanets, Zhenjian Lu, Igor C. Oliveira}
\date{doi:10.4230/LIPIcs.CCC.2022.16; CC BY 4.0}
\begin{document}
\maketitle
\noindent\textit{Editorial scope.} Counted original Theorem43 once: each promised polynomial-hierarchy level has a randomized subexponential simulation under the exact stated higher-level randomized average-case easiness assumption. Complete Theorem40 linear witness compression and all source-local pKt/rKt/PRG support included, not separately counted. No unconditional simulation, deterministic algorithm or collapse of PH to BPP claimed.\par\medskip
\noindent\textit{Allowed standard prerequisites (editorial).}\begin{itemize}
\item Standard polynomial hierarchy/promise classes and Lautemann simulation; imported DPG reconstruction [11] and pseudodeterministic PRG hardness ingredients [19,39,35]. Exact imported PRG/circuit-hardness statements retained. All source-local padding, reconstruction, symmetry, oracle elimination, witness compression and hierarchy induction proofs included.
\end{itemize}\medskip
% Faithful official CCC2022 article16 excerpts: prerequisites and complete pKt proof chain for the UP main theorem.
\section*{2 Preliminaries (necessary author context)}
\subsection*{2.1 Notation}
We say that a function $p:\mathbb N\to\mathbb N$ is non-decreasing if $p(\ell)\geq\ell$ for every integer $\ell\geq0$. Throughout the paper, we will rely on constructions whose associated complexity parameters (e.g., running time) can be bounded by some polynomial $p$. We will implicitly assume that the polynomials provided in these results are non decreasing in order to avoid including this condition in every statement. This can be done without loss of generality by the monotonicity of our bounds.
For a probability distribution $D_n$, we use $\operatorname{supp}(D_n)$ to denote its support. Given an element $x$, we let $D_n(x)$ denote the probability of $x$ under $D_n$. We denote by $U_n$ the uniform distribution over $n$-bit strings.
Let $D_1$ and $D_2$ be probability distributions supported over a set $X$, $A:X\to\{0,1\}$, and $\varepsilon\geq0$. We say that $A$ $\varepsilon$-distinguishes $D_1$ and $D_2$ if
\[\left|\Pr_{x\sim D_1}[A(x)=1]-\Pr_{x\sim D_2}[A(x)=1]\right|\geq\varepsilon.\]
If $A$ is a randomized algorithm, each probability in the expression above also takes into account the internal randomness of $A$, which will be denoted by $r_A$.
\subsection*{2.2 Average-case complexity}
Recall that a pair $(L,\mathcal D)$ is a distributional problem if $L\subseteq\{0,1\}^*$ and $\mathcal D=\{D_n\}_{n\geq1}$ is an ensemble of probability distributions, where each $D_n$ is supported over $\{0,1\}^*$.
Let $\mathcal D=\{D_n\}_{n\geq1}$ be an ensemble of distributions. We say that $\mathcal D\in\mathrm{PSamp}$ if there is a randomized polytime algorithm $A$ such that, for every $n\geq1$, $A(1^n)$ is distributed according to $D_n$. More generally, we use $D_n\in\mathrm{Samp}[T(n)]/a(n)$ to denote that $D_n$ can be sampled by an algorithm that runs in time $T(n)$ and has advice complexity $a(n)$.
We let $\mathrm{DistNP}$ denote the set of distributional problems $(L,\mathcal D)$ with $L\in\mathrm{NP}$ and $\mathcal D\in\mathrm{PSamp}$.
\paragraph{Definition 4 ($\mathrm{Avg}_\delta\mathrm{BPP}$ [3]).} Let $(L,\mathcal D)$ be a distributional problem, and $\delta:\mathbb N\to[0,1]$. We say that $(L,\mathcal D)\in\mathrm{Avg}_\delta\mathrm{BPP}$ if there is randomized polytime algorithm $A$ such that
\begin{enumerate}
\item for every $n>0$, and every $x\in\operatorname{supp}(D_n)$,
\[\Pr_{r_A}[A(x;n)\notin\{L(x),\bot\}]\leq1/4,\]
\item for every $n>0$,
\[\Pr_{x\sim D_n}\left[\Pr_{r_A}[A(x;n)=\bot]\geq1/4\right]\leq\delta(n),\]
\end{enumerate}
where $r_A$ denotes the internal randomness of $A$. Such an algorithm $A$ is called an randomized errorless heuristic for $(L,\mathcal D)$ with failure probability at most $\delta$.
\paragraph{Definition 5 (AvgBPP [3]).}\footnote{Some authors define $\mathrm{AvgBPP}$ as $\bigcap_{c>0}\mathrm{Avg}_{n^{-c}}\mathrm{BPP}$. All our results also hold with respect to that definition.} We say that $(L,\mathcal D)$ is in $\mathrm{AvgBPP}$ if there exist a randomized algorithm $A$ and a polynomial $p$ such that
\begin{enumerate}
\item for every $n,\delta>0$, and every $x\in\operatorname{supp}(D_n)$,
\[\Pr_{r_A}[A(x;n,\delta)\notin\{L(x),\bot\}]\leq1/4,\]
\item for every $n,\delta>0$,
\[\Pr_{x\sim D_n}\left[\Pr_{r_A}[A(x;n,\delta)=\bot]\geq1/4\right]\leq\delta(n),\]
\item for every $n,\delta>0$, and every $x\in\operatorname{supp}(D_n)$, $A(x;n,\delta)$ runs in time at most $p(n/\delta)$.
\end{enumerate}
Such an algorithm $A$ is called a randomized errorless heuristic scheme for $(L,\mathcal D)$.
\paragraph{Lemma 6.} Let $(L,\mathcal D)$ be a distributional problem in $\mathrm{Avg}_{n^{-1}}\mathrm{BPP}$. There exists a randomized polynomial-time algorithm $B$ such that
\begin{enumerate}
\item for every $n>0$, and every $x\in\operatorname{supp}(D_n)$, $B(x;n)\in\{0,1\}$,
\item for every $n>0$, and every $x\in\operatorname{supp}(D_n)$, if $x\in L$, then $\Pr_{r_B}[B(x;n)=1]\geq1-1/n$, and
\item for every $n>0$, $\Pr_{x\sim D_n,r_B}[B(x;n)=L(x)]\geq1-3/n$.
\end{enumerate}
\paragraph{Proof.} Let $A$ be an $\mathrm{Avg}_{n^{-1}}\mathrm{BPP}$ algorithm for $(L,\mathcal D)$ as in Definition 4. Let $A'$ be the algorithm obtained, using standard amplification techniques,\footnote{In more detail, $A'$ can be obtained from $A$ by running it $O(n)$ times and selecting the most common output.} such that for all $x\in\operatorname{supp}(D_n)$,
\[\Pr_{r_{A'}}[A'(x;n)\notin\{L(x),\bot\}]\leq1/n\tag{1}\]
and
\[\Pr_{x\sim D_n}\left[\Pr_{r_{A'}}[A'(x;n)=\bot]\geq1/n\right]\leq1/n.\tag{2}\]
Let $B$ be the randomized algorithm which, on input $x\in\operatorname{supp}(D_n)$, simulates $A'(x;n)$ and outputs $1$ if $A'(x;n)\in\{1,\bot\}$ or $0$ if $A'(x;n)=0$. By Equation (1), if $x\in L$, then
\[\Pr_{r_B}[B(x;n)=0]=\Pr_{r_{A'}}[A'(x;n)\notin\{L(x),\bot\}]\leq1/n.\]
Let $\mathrm{bad}:=\{x\in\operatorname{supp}(D_n)\mid\Pr_{r_{A'}}[A'(x;n)=\bot]\geq1/n\}$. By Equation (2), at most a $(1/n)$-measure of $x\in\operatorname{supp}(D_n)$ are in bad. Thus,
\begin{align*}
\Pr_{x\sim D_n,r_B}[B(x;n)\neq L(x)]
&\leq\Pr_{x,r_{A'}}[A'(x;n)=\bot]+\Pr_{x,r_{A'}}[A'(x;n)\notin\{L(x),\bot\}]\\
&\leq\Pr_x[x\in\mathrm{bad}]+\Pr_{x,r_{A'}}[A'(x;n)=\bot\mid x\notin\mathrm{bad}]\\
&\qquad+\Pr_{x,r_{A'}}[A'(x;n)\notin\{L(x),\bot\}]\\
&\leq1/n+1/n+1/n=3/n.
\end{align*}
This completes the proof.\hfill$\triangleleft$
\subsection*{2.4 Direct product generators (DPGs)}
\paragraph{Definition 13 (Direct Product Generator (DPG) [11, Definiton 3.10]).} For $k,n\in\mathbb N$, we define the $k$-wise direct product generator to be the function
\[\mathrm{DP}_k:\{0,1\}^n\times\{0,1\}^{nk}\to\{0,1\}^{nk+k}\]
such that
\[\mathrm{DP}_k(x;z^1,\ldots,z^n):=(z^1,\ldots,z^k,x\cdot z^1,\ldots,x\cdot z^k).\]
\paragraph{Lemma 14 (DPG Reconstruction [11, Lemma 3.14]).} For any $n,k\in\mathbb N$ with $k\leq2n$, and $\varepsilon>0$, there exists a pair of algorithms $A$ and $\mathrm{Recon}^{(-)}$ such that
\begin{itemize}
\item $\mathrm{Recon}^D$ takes oracle access to an oracle $D:\{0,1\}^{nk+k}\to\{0,1\}$.
\item $A:\{0,1\}^n\times\{0,1\}^r\to\{0,1\}^k$ is called the advice function and is computable in time $\operatorname{poly}(n/\varepsilon)$.
\item $\mathrm{Recon}^{(-)}:\{0,1\}^k\times\{0,1\}^r\to\{0,1\}^n$ is called a reconstruction procedure and is computable in time $\operatorname{poly}(n/\varepsilon)$.\footnote{We note that if the oracle is a uniform algorithm that runs in time $t(m)$ on inputs of length $m$, the reconstruction procedure can be computed in time $t(nk+k)\cdot\operatorname{poly}(n/\varepsilon)$.}
\item The randomness complexity $r$ is at most $\operatorname{poly}(n/\varepsilon)$.
\item For any string $x\in\{0,1\}^n$ and any function $D$ that $\varepsilon$-distinguishes $\mathrm{DP}_k(x;U_{nk})$ from $U_{nk+k}$, it holds that
\[\Pr_{w\sim\{0,1\}^r}[\mathrm{Recon}^D(A(x,w),w)=x]\geq1/\operatorname{poly}(n/\varepsilon).\]
\end{itemize}
\subsection*{2.7 Encoding tuples (author context)}
We will also need to consider distributions that generate input instances that are tuples, while the only known information to the distributions is the length of the instances. To do this, we will encode the format as the length of the instances. More specifically, we use the following way to encode tuples [3].
\paragraph{Proposition 16 ([3]; see also [6, Proposition 3.1]).} There is a pair of polytime encoding-decoding algorithms $\mathrm{Enc}:\mathbb N^*\to\mathbb N$ and $\mathrm{Dec}:\mathbb N\to\mathbb N^*\cup\{\bot\}$ such that
\begin{enumerate}
\item For $k\in\mathbb N$ and $k$ integers $n_1,n_2,\ldots,n_k$, we write $\langle n_1,n_2,\ldots,n_k\rangle:=\mathrm{Enc}(n_1,n_2,\ldots,n_k)$, and we have
\begin{enumerate}
\item $\sum_{i=1}^k(2\cdot(\lfloor\log n_i\rfloor+2)+n_i)\leq\langle n_1,\ldots,n_k\rangle\leq O_k(\prod_{i=1}^k(n_i\cdot\log^2n_i))$.
\item $\mathrm{Dec}(\langle n_1,n_2,\ldots,n_k\rangle)=n_1,\ldots,n_k$. Here, Dec takes the binary representation of an integer and outputs $k$ integers which are also represented in binary.
\end{enumerate}
\item $\mathrm{Dec}(u)=\bot$ if $u\neq\langle\vec n\rangle$ for any $\vec n\in\mathbb N^*$.
\end{enumerate}
\paragraph{Proof.} Given $k$ integers $n_1,n_2,\ldots,n_k$, the encoding algorithm Enc first obtains the binary string
\[z:=\alpha_1\circ01\circ\alpha_2\circ01\circ\alpha_k\circ01\circ\operatorname{bin}(n_1)\circ\operatorname{bin}(n_2)\circ\cdots\circ\operatorname{bin}(n_k),\]
where for each $i\in[k]$, $\alpha_i$ is obtained from $\operatorname{bin}(|\operatorname{bin}(n_i)|)$ with every bit doubled. Then Enc outputs the integer $n$ whose binary representation is $z$. It is easy to see that given the binary representation of $n$, which is simply the string $z$, one can recover the $n_1,n_2,\ldots,n_k$ in time polynomial in $|z|$. Also, item (a) can be verified by a simple calculation.\hfill$\triangleleft$
\section*{3 Probabilistic Time-Bounded Kolmogorov Complexity (author excerpts)}
In this section, we formally define probabilistic Kolmogorov complexity and prove some useful properties of this notion, which will be used later in proving our main results.
\paragraph{Definition 17 ($\mathrm{pK}^t$).} For strings $x,y\in\{0,1\}^*$, a time bound $t\in\mathbb N$, an oracle $A$, and $\delta\in[0,1]$, the $\delta$-probabilistic $A$-oracle $t$-bounded Kolmogorov complexity of $x$ given $y$ is defined as
\[\mathrm{pK}_\delta^{t,A}(x\mid y):=\min\left\{k\ \middle|\ \Pr_{w\sim\{0,1\}^t}[\exists M\in\{0,1\}^k,M^A(w,y)\text{ outputs }x\text{ within }t\text{ steps}]\geq\delta\right\},\]
where $M$ is a RAM-machine.\footnote{We use the RAM model in this definition for convenience: it simplifies the time bound estimates for a machine $M$ that takes as input both randomness $w$ and an ``auxiliary'' string $y$, and allows both $w$ and $y$ to be as long as the running time of the RAM-machine $M$; cf. [29] for a similar definition in the case of deterministic time-bounded conditional Kolmogorov complexity. Alternatively, we could use a TM model, with inputs $w$ and $y$ presented on two different tapes. However, we prefer the RAM model as it's also useful in the context of fine-grained worst-case to average-case reductions [6].} We omit the subscript $\delta$ if $\delta=2/3$, omit the superscript $A$ if $A=\emptyset$, omit ``$\mid y$'' if $y$ is the empty string, and omit the superscript $t$ if $t=\infty$.
For simplicity, in the rest of this section we only consider $\mathrm{pK}^t$ without oracles. It is easy to see that all the results also hold with any oracle.
\paragraph{Lemma 21 (Success Amplification).} For any string $x\in\{0,1\}^n$, time bound $t\in\mathbb N$, and $0\leq\alpha<\beta\leq1$, we have
\[\mathrm{pK}_\beta^{O(qt/\alpha)}(x)\leq\mathrm{pK}_\alpha^t+O(\log(q/\alpha)),\]
where $q=\ln(1/(1-\beta))$.
\paragraph{Proof.} Suppose $\mathrm{pK}_\alpha^t(x)=k$. Then by definition, we have
\[\Pr_{w\sim\{0,1\}^t}[\exists M_w\in\{0,1\}^k\text{ such that }M_w(w)\text{ outputs }x\text{ within }t\text{ steps}]\geq\alpha.\]
Call such a $w$ good. After sampling $\ell$ independent $w^{(1)},w^{(2)},\ldots,w^{(\ell)}\in\{0,1\}^t$, the probability that no $w^{(i)}$ is good is at most $(1-\alpha)^\ell\leq e^{-\alpha\ell}$, which can be made at most $1-\beta$ by choosing $\ell=q/\alpha$. It follows that, with probability at least $\beta$ over a random $w=(w^{(1)},w^{(2)},\ldots,w^{(q/\alpha)})\in\{0,1\}^{qt/\alpha}$, there exists an index $1\leq i\leq q/\alpha$ (described with at most $\lceil\log(q/\alpha)\rceil$ bits) and a program $M_{w^{(i)}}\in\{0,1\}^k$ such that $M_{w^{(i)}}(w^{(i)})$ outputs $x$ within $O(qt/\alpha)$ steps. Hence, $\mathrm{pK}_\beta^{O(qt/\alpha)}(x)\leq k+O(\log(q/\alpha))$.\hfill$\triangleleft$
Note that the parameter $\alpha$ affects the time complexity in the resulting $\mathrm{pK}_\beta^{O(qt/\alpha)}$ bound. Lemma 21 is sufficient in applications where $\alpha\geq1/\operatorname{poly}(n)$.
\subsection*{3.2 Bounding $\mathrm{pK}^t$ via DPG reconstruction (selected original case)}
The following is a key lemma for $\mathrm{pK}^t$ that we will use in proving our main results related to randomized average-case complexity.
\paragraph{Lemma 22 ($\mathrm{pK}^t$ Reconstruction Lemma).} For $\varepsilon>0$, $x\in\{0,1\}^n$, $s\in\mathbb N$, and $k\in\mathbb N$ satisfying $k\leq2n$, let $D$ be a randomized algorithm that takes an advice string $\beta$ and runs in time $t_D$ such that $D$ $\varepsilon$-distinguishes $\mathrm{DP}_k(x;U_{nk})$ from $U_{nk+k}$. Then there is a polynomial $p_{\mathrm{DP}}$ such that
\[\mathrm{pK}^{\widetilde O(t_D)\cdot p_{\mathrm{DP}}(n/\varepsilon)}(x\mid\beta)\leq k+\log p_{\mathrm{DP}}(nt_D/\varepsilon).\]
\paragraph{Proof.} We will view the distinguisher $D$ as a deterministic algorithm that takes a string of $t_D$ bits, denoted by $r_D$, as its internal randomness. By our assumption (and dropping the absolute value sign without loss of generality), we have
\[\mathbb E_{y\sim\{0,1\}^{nk},r_D\sim\{0,1\}^{t_D}}[D(\mathrm{DP}_k(x,y);r_D)]-\mathbb E_{z\sim\{0,1\}^{nk+k},r_D\sim\{0,1\}^{t_D}}[D(z;r_D)]\geq\varepsilon.\]
Then by an averaging argument, we have
\[\Pr_{r_D}\left[\mathbb E_y[D(\mathrm{DP}_k(x,y);r_D)]-\mathbb E_z[D(z;r_D)]\geq\varepsilon/2\right]\geq\varepsilon/2.\tag{4}\]
In other words, with probability at least $\varepsilon/2$ over the internal randomness $r_D$, $D(-,r_D)$ is an $(\varepsilon/2)$-distinguisher for $\mathrm{DP}_k(x;U_{nk})$ and $U_{nk+k}$. Let us say that $r_D$ is good if this is true.
Let $\mathrm{Recon}^{(-)}:\{0,1\}^k\times\{0,1\}^r\to\{0,1\}^n$ be the reconstruction procedure in Lemma 14 that works for distinguishing parameter $\varepsilon/2$, where the randomness complexity $r$ is $\operatorname{poly}(n/\varepsilon)$. We have
\begin{align*}
&\Pr_{w\sim\{0,1\}^r,r_D\sim\{0,1\}^{t_D}}[\exists\alpha\in\{0,1\}^k\text{ such that }\mathrm{Recon}^{D(-,r_D)}(\alpha,w)=x]\\
&\geq\Pr_{w,r_D}[\exists\alpha\in\{0,1\}^k\text{ such that }\mathrm{Recon}^{D(-,r_D)}(\alpha,w)=x\mid r_D\text{ is good}]\cdot\Pr_{r_D}[r_D\text{ is good}]\\
&\geq\frac1{\operatorname{poly}(n/\varepsilon)}\cdot(\varepsilon/2).\qquad\text{(Lemma 14 and Equation (4))}
\end{align*}
The above implies that for at least $1/\operatorname{poly}(n/\varepsilon)$ fraction of the randomness $(w,r_D)$, there exists a program $M$, which takes $k$ bits for some $\alpha$ (that can depend on the randomness) and the advice of $\beta$ used by $D$, such that $M$ simulates $\mathrm{Recon}^{D(-,r_D)}$ on $(\alpha,w)$ and outputs $x$. Note that since $\mathrm{Recon}^D$ runs in time $t_D\cdot\operatorname{poly}(n/\varepsilon)$, our program $M$ can be made to run in time $T:=\widetilde O(t_D)\cdot\operatorname{poly}(n/\varepsilon)$. Therefore, we have
\[\mathrm{pK}_{1/\operatorname{poly}(n/\varepsilon)}^T(x\mid\beta)\leq k+O(\log(nt_D/\varepsilon)).\]
Then the lemma follows easily from Lemma 21.\hfill$\triangleleft$

\subsection*{3.4 Symmetry of information (original Item1 and its full proof)}
\paragraph{Lemma 26 (Symmetry of information for $\mathrm{pK}^t$ and $\mathrm{rK}^t$ under the average-case easiness of NP), Item 1.} If $\mathrm{DistNP}\subseteq\mathrm{AvgBPP}$, then there exist polynomials $p$ and $p_0$ such that for all sufficiently large $n,m\in\mathbb N$, all $t\geq p_0(n,m)$, and all $0\leq\tau\leq t$ the following holds: for every $x\in\{0,1\}^n$ and every family $\{y_r\in\{0,1\}^m\mid r\in\{0,1\}^\tau\}$, with probability at least $9/10$ over $r\in\{0,1\}^\tau$,
\[\mathrm{pK}^t(x,y_r\mid r)>\mathrm{pK}^{p(t)}(x\mid r)+\mathrm{pK}^{p(t)}(y_r\mid x,r)-\log p(t).\]
Moreover, the special case without conditioning on $r$ also holds for both items above, i.e., when $r$ is the empty string (for $\tau=0$) and $y$ is an arbitrary single string.
\paragraph{Proof of Item 1.} Define a language
\[L:=\{(u,v,w,w',1^s)\mid\exists M\in\{0,1\}^s,M(w,w')\text{ prints }uv\text{ in }|w|\text{ steps, and }s=|u|+|v|-10\}.\]
Note that $L\in\mathrm{NP}$. Define a distribution family $\mathcal D=\{D_{\langle nk+k,mk'+k',2t,\tau,s\rangle}\}$ as follows: sample $u\sim U_{nk+k}$, $v\sim U_{mk'+k'}$, $w\sim U_{2t}$, and $w'\sim U_\tau$, and then output $(u,v,w,w',1^s)$. Note that the ensemble $\mathcal D$ is in PSamp under an appropriate encoding of its defining tuple (for simplicity, we omit the padding in the output of $\mathcal D$; see Section 2.1). Using the assumption that $\mathrm{DistNP}\subseteq\mathrm{AvgBPP}$, it follows that $(L,\mathcal D)\in\mathrm{AvgBPP}$. Let $B$ be a randomized algorithm for $(L,\mathcal D)$ as in Lemma 6.
Let $x\in\{0,1\}^n$ and $\{y_r\}_{r}\subseteq\{0,1\}^m$ be given (for sufficiently large $n$ and $m$). Let $k,k'\in\mathbb N$ be arbitrary parameters such that $k\leq n$ and $k'\leq m$. Observe that there exists a polynomial $p_0$ such that for any $t\geq p_0(n,m)$, some constant $d$, and all $z\in\{0,1\}^{nk}$, $z'\in\{0,1\}^{mk'}$, and $r\in\{0,1\}^\tau$,
\[\mathrm{pK}^{2t}(\mathrm{DP}_k(x;z),\mathrm{DP}_{k'}(y_r;z')\mid r)\leq\mathrm{pK}^t(x,y_r\mid r)+|z|+|z'|+d\log t.\tag{6}\]
Here $p_0(n,m)$ reflects the time required to deterministically compute $(\mathrm{DP}_k(x;z),\mathrm{DP}_{k'}(y_r;z'))$ given $x,y_r,z,z'$, and $d\log t$ bits of information to delineate $x$ from $y_r$.
Let $t\geq p_0(n,m)$. By a counting argument, for independently random $u,v,w$, and $w'$,
\[\Pr_{u,v,w,w'}[(u,v,w,w',1^s)\in L]\leq\frac{2^s\cdot2^{|w|+|w'|}}{2^{|u|+|v|+|w|+|w'|}}=2^{s-|u|-|v|}=2^{-10},\]
where the last equality is by the definition of $L$ requiring that $s=|u|+|v|-10$. Then
\begin{align*}
\Pr_{u,v,w,w',r_B}[B(u,v,w,w',1^s)=1]
&\leq\Pr_{u,v,w,w'}[(u,v,w,w',1^s)\in L]\\
&\quad+\Pr_{u,v,w,w',r_B}[B(u,v,w,w',1^s)\neq L(u,v,w,w',1^s)]\\
&\leq2^{-10}+(3/n).\tag{7}
\end{align*}
By Markov's inequality, for at least $9/10$ fraction of random strings $r\in\{0,1\}^\tau$, we get
\[\Pr_{u,v,w,r_B}[B(u,v,w,r,1^s)=1]\leq10\cdot(2^{-10}+(3/n))\leq1/10.\tag{8}\]
Fix any such $r$ so that Equation (8) holds. This also fixes the string $y_r$.
Next we show, using a hybrid argument, that $B(-,U_{2t},r,1^s)$ cannot distinguish between the uniform distribution and the distribution $(\mathrm{DP}_k(x;z),\mathrm{DP}_{k'}(y_r;z'))$, for random independent $z,z'$, where $k\approx\mathrm{pK}^{p_{\mathrm{DP}}(t)}(x\mid r)$ and $k'\approx\mathrm{pK}^{p_{\mathrm{DP}}(t)}(y_r\mid x,r)$. This will imply that
\[\Pr_w[(\mathrm{DP}_k(x;z),\mathrm{DP}_{k'}(y_r;z'),w,r,1^s)\in L]<2/3\]
for some $z,z'$, yielding the desired lower bound on $\mathrm{pK}^{2t}(\mathrm{DP}_k(x;z),\mathrm{DP}_{k'}(y_r;z')\mid r)$ via our choice of $s$ and the definition of $L$. We give the details of the hybrid argument next.
First, toward a contradiction, suppose
\[\Pr_{z,v,w,r_B}[B(\mathrm{DP}_k(x;z),v,w,r,1^s)=1]>1/2.\]
In this case, comparing with Equation (8), we get a randomized distinguisher (with advice) for $\mathrm{DP}_k(x;U_{nk})$, defined by sampling $v\sim U_{mk'+k'}$ and $w\sim U_{2t}$, and outputting $B(-,v,w,r,1^s)$. By Lemma 22,
\[\mathrm{pK}^{q(t)}(x\mid r)\leq k+\log q(t)\tag{9}\]
for some polynomial $q$ such that $q(t)\geq p_{\mathrm{DP}}(t_B\cdot3\cdot n)$ whenever $t\geq n+m$ and $s\leq n^3$, where $p_{\mathrm{DP}}$ is the polynomial from Lemma 22 and $t_B$ denotes the time required to compute $B(-,v,w,r,1^s)$.
Define $k:=\mathrm{pK}^{q(t)}(x\mid r)-\log q(t)-1$ so that Equation (9) does not hold. Assume for now that $k>0$. Hence,
\[\Pr_{z,v,w,r_B}[B(\mathrm{DP}_k(x;z),v,w,r,1^s)=1]\leq1/2.\tag{10}\]
Again, toward a contradiction, suppose that for all $z,z'$,
\[\Pr_w[(\mathrm{DP}_k(x;z),\mathrm{DP}_{k'}(y_r;z'),w,r,1^s)\in L]\geq2/3.\]
By definition of $B$, this implies that
\[\Pr_{z,z',w,r_B}[B(\mathrm{DP}_k(x;z),\mathrm{DP}_{k'}(y_r;z'),w,r,1^s)=1]\geq(2/3)(1-1/n)>5/8.\]
In this case, comparing with Equation (10), we get a randomized distinguisher (with advice) $B'$ for $\mathrm{DP}_{k'}(y_r;U_{mk'})$, defined by sampling $z\sim U_{nk}$, $w\sim U_{2t}$, and outputting $B(\mathrm{DP}_k(x;z),-,w,r,1^s)$. By Lemma 22,
\[\mathrm{pK}^{q'(t)}(y_r\mid x,r)\leq k'+\log q'(t)\tag{11}\]
for some polynomial $q'$ with $q'(t)\geq p_{\mathrm{DP}}(t_{B'}\cdot8\cdot m)$ whenever $t\geq n+m$ and $s\leq n^3$, where $t_{B'}$ denotes the time required to compute $B'$.
We now choose $k':=\mathrm{pK}^{q'(t)}(y_r\mid x,r)-\log q'(t)-1$ so that Equation (11) does not hold. Assume for now that $k'>0$. Hence, there exist $z$ and $z'$ such that
\begin{align*}
&\Pr_w[(\mathrm{DP}_k(x;z),\mathrm{DP}_{k'}(y_r;z'),w,r,1^s)\in L]\\
&=\Pr_w[\exists M\in\{0,1\}^s,M(w,r)\text{ outputs }(\mathrm{DP}_k(x;z),\mathrm{DP}_{k'}(y_r;z'))\text{ within }|w|\text{ steps}]<2/3,
\end{align*}
which implies that
\[\mathrm{pK}^{2t}(\mathrm{DP}_k(x;z),\mathrm{DP}_{k'}(y_r;z')\mid r)>s.\]
For this choice of $z,z'$, by definition of $s$ we get that $s=|u|+|v|-10=|z|+k+|z'|+k'-10$, and so
\[\mathrm{pK}^{2t}(\mathrm{DP}_k(x;z),\mathrm{DP}_{k'}(y_r;z')\mid r)>|z|+k+|z'|+k'-10.\]
Combining this inequality with Equation (6), we get
\begin{align*}
\mathrm{pK}^t(x,y_r\mid r)
&\geq\mathrm{pK}^{2t}(\mathrm{DP}_k(x;z),\mathrm{DP}_{k'}(y_r;z')\mid r)-|z|-|z'|-d\log t\\
&>k+k'-d\log t-10.\tag{12}
\end{align*}
The definitions of $k$ and $k'$ then imply that
\begin{multline*}
\mathrm{pK}^t(x,y_r\mid r)>\mathrm{pK}^{q(t)}(x\mid r)+\mathrm{pK}^{q'(t)}(y_r\mid x,r)\\
-\log q(t)-\log q'(t)-d\log t-12.
\end{multline*}
For the polynomial $p(t):=(t^{d+1})\cdot q(t)\cdot q'(t)$ and for every $t\geq p_0(|x|,|y|)$, we get
\[\mathrm{pK}^t(x,y_r\mid r)>\mathrm{pK}^{p(t)}(x\mid r)+\mathrm{pK}^{p(t)}(y_r\mid x,r)-\log p(t),\]
as desired.
Finally, consider the case that $k\leq0$ or $k'\leq0$. If $k\leq0$, then $\mathrm{pK}^{q(t)}(x\mid r)\leq\log q(t)+1$, implying that $\mathrm{pK}^{p(t)}(x\mid r)<\log p(t)$. But then the lemma simply follows from the fact that $\mathrm{pK}^t(x,y_r\mid r)\geq\mathrm{pK}^{p(t)}(y_r\mid x,r)$. Similarly, if $k'\leq0$, then $\mathrm{pK}^{p(t)}(y_r\mid x,r)<\log p(t)$, and the lemma follows from the fact that $\mathrm{pK}^t(x,y_r\mid r)\geq\mathrm{pK}^{p(t)}(x\mid r)$.\hfill$\triangleleft$
\subsection*{3.7 Probabilistic computational depth (original Item1 and its full proof)}
The next lemma is a straightforward adaptation of an argument from [11]. We provide a proof for completeness.
\paragraph{Lemma 31 (Computational Depth Upper Bound), Item 1.} For all $x\in\{0,1\}^n$ and $r\in\{0,1\}^*$, and for all non-decreasing polynomials $p_0$ and $p$, if $n$ is large enough then there exists a time bound $t_1$ such that $p_0(n)\leq t_1\leq2^{n/\log n}$ and
\[\mathrm{pK}^{t_1}(x\mid r)-\mathrm{pK}^{p(t_1)}(x\mid r)\leq O(n/\log n).\]
\paragraph{Proof.} We will start by proving Item 1. Given $x\in\{0,1\}^n$ and polynomials $p_0$ and $p$, define the polynomial $\tau:=p\circ p_0$. For an integer $I\geq1$, consider the following telescoping sum:
\begin{align*}
\mathrm{pK}^{\tau(n)}(x\mid r)-\mathrm{pK}^{\tau^{I+1}(n)}(x\mid r)
&=(\mathrm{pK}^{\tau(n)}(x\mid r)-\mathrm{pK}^{\tau^2(n)}(x\mid r))\\
&\quad+(\mathrm{pK}^{\tau^2(n)}(x\mid r)-\mathrm{pK}^{\tau^3(n)}(x\mid r))\\
&\quad+\cdots+(\mathrm{pK}^{\tau^I(n)}(x\mid r)-\mathrm{pK}^{\tau^{I+1}(n)}(x\mid r)),
\end{align*}
where $\tau^i(-)$ denotes the composition of $\tau$ with itself $i$ times. For any choice of $x,p_0$, and $p$ as in the statement of the lemma, $\mathrm{pK}^{\tau(n)}(x\mid r)\leq n+d$, for some universal constant $d\geq0$; hence, the above sum is at most $n+d$. By averaging, there is some index $i_0\in[I]$ such that
\[\mathrm{pK}^{\tau^{i_0}(n)}(x\mid r)-\mathrm{pK}^{\tau^{i_0+1}(n)}(x\mid r)\leq(n+d)/I.\]
For this $i_0$, define $t_1:=\tau^{i_0}(n)$. Note that $t_1\geq\tau(n)\geq p_0(n)$, since $i_0\geq1$ and $p(\ell)\geq\ell$ for every input $\ell$. Letting $c\in\mathbb N$ be such that $\tau(n)\leq n^c$ for sufficiently large $n$, define $I:=\log_c(n/(\log_2n)^2)$. Then $t_1\leq n^{c^I}=2^{n/\log n}$. Moreover,
\begin{align*}
\mathrm{pK}^{t_1}(x\mid r)-\mathrm{pK}^{p(t_1)}(x\mid r)&\leq\mathrm{pK}^{t_1}(x\mid r)-\mathrm{pK}^{\tau(t_1)}(x\mid r)\\
&\leq O(n/\log n),
\end{align*}
as desired.

% Faithful official CCC2022 article16 original pKt cases of Lemmas33–34 and Theorem36.
\section*{4 Probabilistic Worst-Case to Average-Case Reductions (AM case)}
\subsection*{4.1 Auxiliary lemmas}
Given a language $L\in\mathrm{NP}$ with a corresponding verifier $V$, and $x\in L$, let $y_x$ be the lexicographically first witness that $x\in L$. It is not hard to compute $y_x$ from $x$ if an NP oracle is available, by performing a standard search-to-decision reduction. For this reason, the time-bounded Kolmogorov complexity of the pair $(x,y_x)$ is essentially that of $x$, in the presence of an NP oracle. We will show, more generally, that the oracle can be eliminated under an appropriate average-case easiness assumption, even for the following randomized versions of $\Sigma_\ell^\mathrm P$, denoted $\mathrm{BP}_\delta\circ\Sigma_\ell^\mathrm P$.
\paragraph{Definition 32 ($\mathrm{BP}_\delta\circ\Sigma_\ell^\mathrm P$).} A language $L$ is in the class $\mathrm{BP}_\delta\circ\Sigma_\ell^\mathrm P$ if there is a $\mathrm{BP}\circ\Sigma_\ell\circ\mathrm P$ formula
\[\varphi(x)=\mathrm{BP}_r\exists y_1\forall y_2\ldots Qy_\ell\ R(x,y_1,y_2,\ldots,y_\ell,r)\]
for a polytime predicate $R$ and all string variables $y_1,\ldots,y_\ell,r$ of length $\operatorname{poly}(|x|)$, such that, for all $x\in\{0,1\}^n$,
\begin{align*}
x\in L&\Longrightarrow\Pr_r[\exists y_1\forall y_2\ldots Qy_\ell\ R(x,y_1,y_2,\ldots,y_\ell,r)]\geq\delta,\text{ and}\\
x\notin L&\Longrightarrow\Pr_r[\forall y_1\exists y_2\ldots\overline Qy_\ell\ \neg R(x,y_1,y_2,\ldots,y_\ell,r)]\geq\delta.
\end{align*}
\paragraph{Lemma 33 (Oracle Elimination), original Item 1.} For any $\ell\geq1$, suppose $\mathrm{Dist}\Sigma_{\ell+1}^{\mathrm P}\subseteq\mathrm{AvgBPP}$. For an arbitrary $L\in\mathrm{BP}_\delta\circ\Sigma_\ell^{\mathrm P}$, let $x\in L_n$ be sufficiently large, let $r\in\{0,1\}^{n^c}$ be any random string under the BP quantifier such that there exists a witness for the left-most $\exists$ quantifier in the definition of $L$, and let $y_{x,r}\in\{0,1\}^{n^c}$ be the lexicographically first such $L$-witness for $x$ and randomness $r$, for some constant $c>0$.
There exist polynomials $q$ and $q_0$ such that, for all sufficiently large $n\in\mathbb N$, all $x\in L_n$, and all $t\geq q_0(n)$, we have with probability at least $\delta-1/10$ over $r\in\{0,1\}^{n^c}$ that an $L$-witness $y_{x,r}$ exists, and
\[\mathrm{pK}^{q(t)}(x,y_{x,r}\mid r)\leq\mathrm{pK}^t(x\mid r)+\log q(t).\]
Moreover, for the case of no BP quantifier, i.e., for $L\in\Sigma_\ell^{\mathrm P}$, we get the same conclusions (with probability $1$) without conditioning on $r$.
\paragraph{Proof of Item 1.} Define
\begin{multline*}
L':=\{(\mathrm{DP}_k(x,y;z),w,w',1^k,1^s,1^n)\mid\exists M\in\{0,1\}^s,\\
M^{\Sigma_\ell^{\mathrm P}}(w,w')\text{ prints }(x,y)\in\{0,1\}^{n+n^c}\text{ within }|w|\text{ steps, where }|w'|=n^c,\text{ and }s=k-10\}.
\end{multline*}
More formally, in this and subsequent proofs, the above notation should be interpreted as follows: a string $(u,w,w',1^k,1^s,1^n)$ belongs to $L'$ iff there exist some $(x,y)\in\{0,1\}^{n+n^c}$ and $M\in\{0,1\}^s$ such that $\mathrm{DP}_k(x,y;z)=u$, where $u=(z,\alpha)$ for some $z\in\{0,1\}^{(n+n^c)k}$ and $\alpha\in\{0,1\}^k$, and $M^{\Sigma_\ell^{\mathrm P}}(w,w')$ outputs $(x,y)$ within $|w|$ steps, for $|w'|=n^c$.
Note that $L'\in\mathrm{NP}^{\Sigma_\ell^{\mathrm P}}$. Define a distribution family $\mathcal D=\{D_{\langle(n+n^c)k+k,2t,n^c,k,s,n\rangle}\}$ as follows: sample $u\sim U_{(n+n^c)k+k}$, $w\sim U_{2t}$, $w'\sim U_{n^c}$, and output $(u,w,w',1^k,1^s,1^n)$, for $s=k-10$. By assumption, $(L',\mathcal D)\in\mathrm{AvgBPP}$. Let $B$ be a randomized algorithm for $(L',\mathcal D)$ as in Lemma 6.
For all (sufficiently large) $n,k$, and $s$, and for independent uniformly random $u,w$, and $w'$,
\begin{align*}
&\Pr_{u,w,w'}[(u,w,w',1^k,1^s,1^n)\in L']\\
&=\Pr_{u,w,w'}[\exists(x,y)\in\{0,1\}^{n+n^c},\exists z\in\{0,1\}^{(n+n^c)k},\exists M\in\{0,1\}^s,\\
&\qquad M^{\Sigma_\ell^{\mathrm P}}(w,w')=(x,y)\wedge u=\mathrm{DP}_k(x,y;z)]\\
&\leq\frac{2^{s+|w|+|w'|+|z|}}{2^{k+|w|+|w'|+|z|}}=2^{s-k}=2^{-10},
\end{align*}
where the inequality is by a union bound and a counting argument, and the last equality by the condition that $s=k-10$ in the definition of $L'$. Hence,
\begin{align*}
\Pr_{u,w,w',r_B}[B(u,w,w',1^k,1^s,1^n)=1]
&\leq\Pr_{u,w,w'}[(u,w,w',1^k,1^s,1^n)\in L']\\
&\qquad+\Pr_{u,w,w',r_B}[B(u,w,w',1^k,1^s,1^n)\neq L'(u,w,w',1^k,1^s,1^n)]\\
&\leq2^{-10}+(3/n).
\end{align*}
By Markov's inequality, for at least $9/10$ fraction of strings $r$,
\[\Pr_{u,w,r_B}[B(u,w,r,1^k,1^s,1^n)=1]\leq10\cdot(2^{-10}+(3/n))\leq1/10.\tag{18}\]
Let $x\in L_n$ be arbitrary. By definition of the $\mathrm{BP}_\delta$ quantifier, for at least $\delta$ fraction of random strings $r\in\{0,1\}^{n^c}$, there exists an $L$-witness, and hence, the lexicographically first $L$-witness $y_{x,r}$. By the above, for at least $\delta-1/10$ of random $r$, we have that both a witness $y_{x,r}$ exists and Equation (18) holds. Fix any such $r$.
Observe that for some polynomial $q_0$ (dependent on $L$) and some constant $d$, for any $t\geq q_0(n)$,
\[\mathrm{pK}^{2t,\Sigma_\ell^{\mathrm P}}(x,y_{x,r}\mid r)\leq\mathrm{pK}^t(x\mid r)+d\log n=:s(x,r).\tag{19}\]
In particular, $q_0(n)$ reflects the time required to deterministically compute $y_{x,r}$, given $x$ and $r$, by a search-to-decision procedure for $L$ using a $\Sigma_\ell^{\mathrm P}$-oracle.
Define $s:=s(x,r)$ (which determines $k=s+10$), and let $t\geq q_0(n)$. By the definitions of $\mathrm{pK}^{2t,\Sigma_\ell^{\mathrm P}}$ and $L'$, for every $z\in\{0,1\}^{(n+n^c)k}$,
\[\Pr_w[(\mathrm{DP}_k(x,y_{x,r};z),w,r,1^k,1^s,1^n)\in L']\geq2/3.\]
Then by the definition of $B$,
\begin{align*}
&\Pr_{z,w,r_B}[B(\mathrm{DP}_k(x,y_{x,r};z),w,r,1^k,1^s,1^n)=1]\\
&\geq\Pr_{z,w}[(\mathrm{DP}_k(x,y_{x,r};z),w,r,1^k,1^s,1^n)\in L']\cdot\Pr_{r_B}[B(\omega)=L'(\omega)\mid\omega\in L']\\
&\geq\frac23(1-(1/n))\geq2/3-o(1).\tag{20}
\end{align*}
Comparing Equation (18) with Equation (20), it is clear that $B(-,U_{2t},r,1^k,1^s,1^n)$ is a randomized algorithm (with advice) that $(1/2)$-distinguishes $\mathrm{DP}_k(x,y_{x,r};U_{(n+n^c)k})$ from uniform. Lemma 22 implies that
\begin{align*}
\mathrm{pK}^{p'(t)}(x,y_{x,r}\mid r,s,n,t)&\leq k+\log p'(t)\tag{21}\\
&=s+10+\log p'(t)\\
&=\mathrm{pK}^t(x\mid r)+d\log n+\log p'(t)+10,
\end{align*}
for some polynomial $p'$ with $p'(t)\geq p_{\mathrm{DP}}(t_B\cdot3\cdot(n+n^c))$, where $p_{\mathrm{DP}}$ is the polynomial from Lemma 22 and $t_B$ denotes the time required to compute $B(-,U_{2t},r,1^k,1^s,1^n)$. As the advice $(s,n,t)$ can be encoded with $\log s+\log n+\log t+O(\log\log n)\leq(3.1)\cdot\log t$ bits, it follows that $\mathrm{pK}^{q(t)}(x,y_{x,r}\mid r)\leq\mathrm{pK}^t(x\mid r)+\log q(t)$ for the polynomial $q(t):=t^{d+4}\cdot p'(t)$.
The ``Moreover'' statement follows the same argument, dropping any mention of $w'$ (and $r$).\hfill$\triangleleft$
\paragraph{Lemma 34 (Witness Compression), original Item 1.} For any $\ell\geq1$, suppose $\mathrm{Dist}\Sigma_{\ell+1}^{\mathrm P}\subseteq\mathrm{AvgBPP}$. For an arbitrary $L\in\mathrm{BP}_\delta\circ\Sigma_\ell^{\mathrm P}$, let $x\in L_n$ be sufficiently large, $r\in\{0,1\}^{n^c}$ be any good random string, and let $y_{x,r}\in\{0,1\}^{n^c}$ be the lexicographically first $L$-witness for $x$ with respect to randomness $r$, for some constant $c>0$.
Let $p_0,p$ and $q_0,q$ be the polynomials from Items 1 of Lemmas 26 and 33 respectively. Then for every $t\geq\max\{p_0(n),q_0(n+n^c)\}$, with probability at least $\delta-1/5$ over $r$, an $L$-witness $y_{x,r}$ exists, and
\[\mathrm{pK}^{p(q(t))}(y_{x,r}\mid x,r)\leq(\mathrm{pK}^t(x\mid r)-\mathrm{pK}^{p(q(t))}(x\mid r))+2\log p(q(t)).\]
Moreover, for $L\in\Sigma_\ell^{\mathrm P}$, we get the same conclusions without conditioning on $r$.
\paragraph{Proof.} We start by proving Item 1. By Item 1 of Lemma 33 and Lemma 26, we get by a union bound that, with probability at least $\delta-2/10$ over $r$,
\[\mathrm{pK}^{q(t)}(x,y_{x,r}\mid r)\leq\mathrm{pK}^t(x\mid r)+\log q(t),\]
and
\[\mathrm{pK}^{q(t)}(x,y_{x,r}\mid r)>\mathrm{pK}^{p(q(t))}(x\mid r)+\mathrm{pK}^{p(q(t))}(y_{x,r}\mid x,r)-\log p(q(t)).\]
Combining the previous two inequalities,
\begin{align*}
\mathrm{pK}^{p(q(t))}(y_{x,r}\mid x,r)
&<\mathrm{pK}^{q(t)}(x,y_{x,r}\mid r)-\mathrm{pK}^{p(q(t))}(x\mid r)+\log p(q(t))\\
&\leq(\mathrm{pK}^t(x\mid r)-\mathrm{pK}^{p(q(t))}(x\mid r))+2\log p(q(t)).
\end{align*}
The ``Moreover'' part follows by applying the ``no random $r$'' versions of Lemma 33 and Lemma 26.
\subsection*{4.2 Case of AM}
\paragraph{Theorem 36.} If $\mathrm{Dist}\Sigma_2^{\mathrm P}\subseteq\mathrm{AvgBPP}$, then $\mathrm{AM}\subseteq\mathrm{BPTIME}[2^{O(n/\log n)}]$.
\paragraph{Proof.} Consider an arbitrary $L\in\mathrm{AM}$ given by a $\mathrm{BP}\circ\Sigma_1^{\mathrm P}$ formula such that, for every $x\in\{0,1\}^n$,
\begin{align*}
x\in L&\Longrightarrow\mathrm{BP}_{(1-2^{-n})}r\ \exists y\ V(x,y,r),\\
x\notin L&\Longrightarrow\mathrm{BP}_{(1-2^{-n})}r\ \forall y\ \neg V(x,y,r),
\end{align*}
for a polytime verifier $V(x,y,r)$, where $|r|,|y|\leq n^c$ for some constant $c>0$, and the notation $\mathrm{BP}_{(\delta)}r\ \varphi(-,r)$, for a formula $\varphi$, means $\Pr_r[\varphi(-,r)]\geq\delta$.
For a given (sufficiently large) $x\in L_n$, and a good random string $r$, let $y_{x,r}$ be the lexicographically first $L$-witness for $x,r$. By Lemma 34, for some constant $a>0$, and for all large enough time bounds $t$, we have, with probability at least $1-2^{-n}-1/5$ over random strings $r$, that
\[\mathrm{pK}^{t^a}(y_{x,r}\mid x,r)\leq(\mathrm{pK}^t(x\mid r)-\mathrm{pK}^{t^a}(x\mid r))+(2a)\log t.\]
By Lemma 31, there exists a time bound $t\leq2^{O(n/\log n)}$ such that $\mathrm{pK}^t(x\mid r)-\mathrm{pK}^{t^a}(x\mid r)\leq bn/\log n$, for some universal constant $b>0$. Hence, for some universal constant $d>0$ (independent of $x,r$),
\[\mathrm{pK}^{2^{dn/\log n}}(y_{x,r}\mid x,r)\leq dn/\log n.\]
A randomized algorithm for $L$ is now obvious:
\begin{itemize}
\item Given $x\in\{0,1\}^n$, sample uniformly random strings $r\in\{0,1\}^{|x|^c}$ and $w\in\{0,1\}^{2^{dn/\log n}}$, and then exhaustively search over all $dn/\log n$-bit machines $M$, running each $M(w,x,r)$ for $2^{dn/\log n}$ steps to produce a candidate witness $y$, accepting if $V(x,y,r)$ accepts.
\end{itemize}
The runtime $2^{O(n/\log n)}$ is clear. For correctness, for every $x\in L_n$, we know there is at least a $(4/5-2^{-n})$ fraction of good random strings $r$ for which a witness $y_{x,r}$ will be found by our algorithm, with probability at least $2/3$ over random $w$'s (by the definition of pK). So our algorithm accepts $x\in L$ with probability at least $(2/3)(4/5)-o(1)>1/2$. On the other hand, for any $x\notin L$, our algorithm may find a witness $y$ and accept $x$ for at most $2^{-n}$ fraction of $r$'s.\hfill$\triangleleft$

% Faithful official CCC2022 article16 additional rKt branches used for the PH main theorem.
\section*{Additional author rKt and pseudodeterministic-PRG context}
We denote by $\mathrm{rK}_\delta^t(x)$ the minimum length of a randomized machine $M$ that outputs $x$ with probability at least $\delta$ when running for at most $t$ steps (see [38,33]). We simply write $\mathrm{rK}^t(x)$ when $\delta=2/3$.
\paragraph{Definition 8 (Pseudodeterministic PRG [39,33]).} Fix an error function $\varepsilon:\mathbb N\to\mathbb R$. A function family $G_n:\{0,1\}^\ell\to\{0,1\}^n$ is an $\varepsilon$-error Pseudorandom Generator (PRG) if, for all sufficiently large $n\in\mathbb N$, the distribution $G(U_{\ell(n)})$ $\varepsilon(n)$-fools circuits of size $n$, using the seed length $\ell=\ell(n)$, i.e., for all circuits $C$ of size $n$ on $n$ inputs,
\[\left|\Pr_{r\sim\{0,1\}^n}[C(r)=1]-\Pr_{\alpha\sim\{0,1\}^{\ell(n)}}[C(G(\alpha))=1]\right|\leq\varepsilon.\]
Such a PRG $G$ is called a pseudodeterministic PRG if there is a randomized algorithm $M_G$ such that, for every $n$ and $\alpha\in\{0,1\}^{\ell(n)}$, $M_G(1^n,\alpha)$ outputs the string $G_n(\alpha)\in\{0,1\}^n$ with probability at least $1-2^{-n}$ over its internal randomness, and $M_G$ runs in time $2^{O(\ell(n))}$.
\paragraph{Lemma 9 ([19,39]).} If BPE contains a language $L$ of circuit complexity $2^{\Omega(n)}$ for almost all input lengths $n$, then there is an $\varepsilon$-error pseudodeterministic PRG of seed length $O(\log(n/\varepsilon))$.
\paragraph{Lemma 10 ([35, Lemma 2]).} $\mathrm E^{\Sigma_2^{\mathrm P}}$ contains a language of maximum circuit complexity (at least $2^n/n$) for almost all input lengths $n$.
Using padding as in [2], we get the following.
\paragraph{Lemma 11.} If $\mathrm{DistP}^{\Sigma_2^{\mathrm P}}\subseteq\mathrm{AvgBPP}$, then $\mathrm E^{\Sigma_2^{\mathrm P}}=\mathrm{BPE}$.
\paragraph{Proof.} The inclusion $\mathrm{BPE}\subseteq\mathrm E^{\Sigma_2^{\mathrm P}}$ follows from the inclusion $\mathrm{BPP}\subseteq\Sigma_2^{\mathrm P}$ [22] by padding. For the other direction, let $L\in\mathrm E^{\Sigma_2^{\mathrm P}}$ be arbitrary. Define its padded version $L_{\mathrm{pad}}=\{(x,1^{2^{2|x|}})\mid x\in L\}$. Note that $L_{\mathrm{pad}}\in\mathrm P^{\Sigma_2^{\mathrm P}}$. Define a family of distributions $D_n=(U_n,1^{2^{2n}})$, for the uniform distribution over $n$-bit strings $U_n$. By assumption, $(L_{\mathrm{pad}},D_n)\in\mathrm{AvgBPP}$, and so, by Lemma 6, Item 3, there is a randomized polytime algorithm $B$ such that $B(x,1^{2^{2|x|}})\neq L(x)$ with probability less than $3\cdot2^{-2n}<1/(3\cdot2^n)$, where the probability is over a uniformly random $x\in\{0,1\}^n$ and the internal randomness of $B$. It follows that, for every $x\in\{0,1\}^n$, $B(x,1^{2^{2|x|}})\neq L(x)$ with probability less than $1/3$ over its internal randomness. Hence, the randomized algorithm $B'(x):=B(x,1^{2^{2|x|}})$ decides if $x\in L$ with probability at least $2/3$, for all sufficiently large inputs $x\in\{0,1\}^n$, yielding $L\in\mathrm{BPE}$.\hfill$\triangleleft$
\paragraph{Corollary 12.} If $\mathrm{DistP}^{\Sigma_2^{\mathrm P}}\subseteq\mathrm{AvgBPP}$, then there is an $\varepsilon$-error pseudodeterministic PRG of seed length $O(\log(n/\varepsilon))$.
\paragraph{Proof.} Immediate by combining Lemmas 9–11.\hfill$\triangleleft$
\paragraph{Lemma 23 ($\mathrm{rK}^t$ Reconstruction Lemma).} Assume the existence of a pseudodeterministic PRG. For $\varepsilon>0$, $x\in\{0,1\}^n$, $s\in\mathbb N$, and $k\in\mathbb N$ satisfying $k\leq2n$, let $D$ be a randomized algorithm that takes an advice $\beta$ and runs in time $t_D$ such that $D$ $\varepsilon$-distinguishes $\mathrm{DP}_k(x;U_{nk})$ from $U_{nk+k}$. Then there is a polynomial $p'_{\mathrm{DP}}$ such that
\[\mathrm{rK}^{p'_{\mathrm{DP}}(nt_D/\varepsilon)}(x\mid\beta)\leq k+\log p'_{\mathrm{DP}}(nt_D/\varepsilon).\]
\paragraph{Proof.} By averaging, arguing as in Lemma 22 up to Equation (4),
\[\Pr_{r_D\sim\{0,1\}^s}[\mathbb E_y[D(\mathrm{DP}_k(x,y),r_D)]-\mathbb E_z[D(z,r_D)]\geq\varepsilon/2]\geq\varepsilon/2.\tag{5}\]
That is, with probability at least $\varepsilon/2$ over its internal randomness $r_D$, $D(-,r_D)$ is an $(\varepsilon/2)$-distinguisher for $\mathrm{DP}_k(x;U_{nk})$ and $U_{nk+k}$. Let us say that $r_D$ is good if this is true.
Let $\mathrm{Recon}^{(-)}:\{0,1\}^k\times\{0,1\}^r\to\{0,1\}^n$ be the reconstruction procedure from Lemma 14 that works for distinguishing parameter $\varepsilon/2$. By definition of Recon, if $r_D$ is good, then
\[\Pr_{w\sim\{0,1\}^r}[\mathrm{Recon}^{D(-,r_D)}(A(x,w),w)=x]\geq1/\operatorname{poly}(n/\varepsilon),\]
and so
\[\Pr_{w,r_D}[\mathrm{Recon}^{D(-,r_D)}(A(x,w),w)=x]\geq\frac\varepsilon2\frac1{\operatorname{poly}(n/\varepsilon)}\stackrel{\mathrm{def}}=\varepsilon'.\]
Observe that by Lemma 14, the condition $\mathrm{Recon}^{D(-,r_D)}(A(x,w),w)=x$ can be checked by a circuit of size $s=\operatorname{poly}(nt_D/\varepsilon)$ given $(r_D,w)$ as input. Consider an assumed pseudodeterministic PRG
\[G:\{0,1\}^\ell\to\{0,1\}^{|r_D|+|w|}\]
that $(\varepsilon'/2)$-fools all circuits of size $s$, where $\ell\in O(\log(s/\varepsilon'))\subseteq O(\log(ns/\varepsilon))$. Let $M_G$ be an algorithm as described in Definition 8, running in time $\operatorname{poly}(s/\varepsilon')\leq\operatorname{poly}(nt_D/\varepsilon)$, such that for all $\sigma\in\{0,1\}^\ell$,
\[\Pr_{r_M\sim\{0,1\}^{\operatorname{poly}(nt_D/\varepsilon)}}[M_G(\sigma,r_M)=G(\sigma)]\geq1-2^{-s},\]
where $r_M$ denotes the internal randomness of $M_G$. By definition of $G$,
\[\Pr_\sigma[\mathrm{Recon}^{D(-,G_0(\sigma))}(A(x,G_1(\sigma)),G_1(\sigma))=x]\geq\varepsilon'/2,\]
where $G_0(\sigma)=G(\sigma)_{1\ldots|r_D|}$ (the $|r_D|$-length prefix) and $G_1(\sigma)=G(\sigma)_{|r_D|+1\ldots|r_D|+|w|}$ (the remaining suffix). So there must exist a seed $\sigma\in\{0,1\}^\ell$ such that
\[\mathrm{Recon}^{D(-,G_0(\sigma))}(A(x,G_1(\sigma)),G_1(\sigma))=x.\]
Let $\sigma$ be a seed with this property, and let $\alpha=A(x,G_1(\sigma))\in\{0,1\}^k$. The definition of $M_G$ implies that $\alpha$ and $\sigma$ are such that
\[\Pr_{r_M}[\mathrm{Recon}^{D(-,M_G^0(\sigma,r_M))}(\alpha,M_G^1(\sigma,r_M))=x]\geq1-2^{-s}>2/3,\]
where $M_G^0(\sigma,r_M)$ denotes the $|r_D|$-length prefix of $M_G(\sigma,r_M)$, and $M_G^1(\sigma,r_M)$ the remaining suffix of $M_G(\sigma,r_M)$. It is easy to verify that the reconstruction procedure $\mathrm{Recon}^{D(-,M_G^0(\sigma,r_M))}(\alpha,M_G(\sigma,-))$ runs in time $\operatorname{poly}(nt_D/\varepsilon)$ overall. By definition of $\mathrm{rK}^t$, for some polynomial $p'_{\mathrm{DP}}$, we get
\begin{align*}
\mathrm{rK}^{p'_{\mathrm{DP}}(nt_D/\varepsilon)}(x\mid\beta)&\leq|\alpha|+|\sigma|+O(\log(ns/\varepsilon))\\
&\leq k+\log p'_{\mathrm{DP}}(nt_D/\varepsilon),
\end{align*}
as required.\hfill$\triangleleft$
\paragraph{Lemma 26, original Item 2.} If $\mathrm{DistP}^{\Sigma_2^{\mathrm P}}\subseteq\mathrm{AvgBPP}$, then there exist polynomials $p'$ and $p'_0$ such that for all sufficiently large $n,m\in\mathbb N$, all $t\geq p'_0(n,m)$, and all $0\leq\tau\leq t$ the following holds: for every $x\in\{0,1\}^n$ and every family $\{y_r\in\{0,1\}^m\mid r\in\{0,1\}^\tau\}$, with probability at least $9/10$ over $r\in\{0,1\}^\tau$,
\[\mathrm{rK}^t(x,y_r\mid r)>\mathrm{rK}^{p'(t)}(x\mid r)+\mathrm{rK}^{p'(t)}(y_r\mid x,r)-\log p'(t).\]
Moreover, the special case without conditioning on $r$ also holds for both items above, i.e., when $r$ is the empty string (for $\tau=0$) and $y$ is an arbitrary single string.
\paragraph{Proof of Item 2.} Argue as in the proof of Item 1, but at Equation (9), apply Lemma 23 instead of Lemma 22 to get
\[\mathrm{rK}^{q(t)}(x\mid r)\leq k+\log q(t)\]
for some polynomial $q$ such that $q(t)\geq p'_{\mathrm{DP}}(t_B\cdot3\cdot n)$, where $p'_{\mathrm{DP}}$ is the polynomial from Lemma 23 and $t_B$ denotes the time required to compute $B(-,v,w,r,1^s)$. We then define $k:=\mathrm{rK}^{q(t)}(x\mid r)-\log q(t)-1$ so the above equation does not hold. Similarly, at Equation (11), apply Lemma 23 to get
\[\mathrm{rK}^{q'(t)}(y_r\mid x,r)\leq k'+\log q'(t)\]
for some polynomial $q'$ with $q'(t)\geq p'_{\mathrm{DP}}(t_{B'}\cdot8\cdot m)$, where $t_{B'}$ denotes the time required to compute the distinguisher $B'$ defined as above. Then define $k':=\mathrm{rK}^{q'(t)}(y_r\mid x,r)-\log q'(t)-1$.
Following the proof of Item 1 up to Equation (12), we get that
\[\mathrm{pK}^t(x,y_r\mid r)>k+k'-d\log t-\log n.\]
The definitions of $k,k'$, and $\mathrm{rK}^t$ then imply that
\begin{multline*}
\mathrm{rK}^t(x,y_r\mid r)>\mathrm{rK}^{q(t)}(x\mid r)+\mathrm{rK}^{q'(t)}(y_r\mid x,r)\\
-\log q(t)-\log q'(t)-d\log t-12,
\end{multline*}
and so
\[\mathrm{rK}^t(x,y_r\mid r)>\mathrm{rK}^{p(t)}(x\mid r)+\mathrm{rK}^{p(t)}(y_r\mid x,r)-\log p(t)\]
for the polynomial $p(t):=(t^{d+1})\cdot q(t)\cdot q'(t)$.\hfill$\triangleleft$
\paragraph{Lemma 33, original Item 2.} Assume in addition that $\mathrm{DistP}^{\Sigma_2^{\mathrm P}}\subseteq\mathrm{AvgBPP}$ (which holds in particular for $\ell\geq2$). Then there exist polynomials $q'$ and $q'_0$ such that, for all sufficiently large $n\in\mathbb N$, all $x\in L_n$, and all $t\geq q'_0(n)$, we have with probability at least $\delta-1/10$ over $r\in\{0,1\}^{n^c}$ that an $L$-witness $y_{x,r}$ exists, and
\[\mathrm{rK}^{q'(t)}(x,y_{x,r}\mid r)\leq\mathrm{rK}^t(x\mid r)+\log q'(t).\]
Moreover, for the case of no BP quantifier, i.e., for $L\in\Sigma_\ell^{\mathrm P}$, we get the same conclusions (with probability $1$) without conditioning on $r$.
\paragraph{Proof of Item 2.} The proof is similar to that of Item 1, except at Equation (19), we observe that for all $t\geq q_0(n)$,
\[\mathrm{pK}^{2t,\Sigma_\ell^{\mathrm P}}(x,y_{x,r}\mid r)\leq\mathrm{pK}^t(x\mid r)+d\log n\leq\mathrm{rK}^t(x\mid r)+d\log n,\]
and then define $s(x,r):=\mathrm{rK}^t(x\mid r)+d\log n$. We then follow the proof of Item 1 up to Equation (21), where we apply Lemma 23 instead of Lemma 22 to obtain
\[\mathrm{rK}^{p''(t)}(x,y_{x,r}\mid r,s,n,t)\leq\mathrm{rK}^t(x\mid r)+d\log n+\log p''(t)+10\]
for some polynomial $p''$ with $p''(t)\geq p'_{\mathrm{DP}}(t_B\cdot3\cdot(n+n^c))$, where $p'_{\mathrm{DP}}$ is the polynomial from Lemma 23 and $t_B$ denotes the time required to compute $B(-,U_{2t},r,1^k,1^s,1^n)$. It follows that $\mathrm{rK}^{q'(t)}(x,y_{x,r}\mid r)\leq\mathrm{rK}^t(x\mid r)+\log q'(t)$ for the polynomial $q'(t):=t^{d+4}\cdot p''(t)$.\hfill$\triangleleft$
\paragraph{Lemma 34, original Item 2.} Assume in addition that $\mathrm{DistP}^{\Sigma_2^{\mathrm P}}\subseteq\mathrm{AvgBPP}$ (which holds in particular when $\ell\geq2$). Let $p'_0,p'$ and $q'_0,q'$ be the polynomials from Items 2 of Lemmas 26 and 33 respectively. Then for every $t\geq\max\{p'_0(n),q'_0(n+n^c)\}$, with probability at least $\delta-1/5$ over $r$, an $L$-witness $y_{x,r}$ exists, and
\[\mathrm{rK}^{p'(q'(t))}(y_{x,r}\mid x,r)\leq(\mathrm{rK}^t(x\mid r)-\mathrm{rK}^{p'(q'(t))}(x\mid r))+2\log p'(q'(t)).\]
Moreover, for $L\in\Sigma_\ell^{\mathrm P}$, we get the same conclusions without conditioning on $r$.
The proof of Item 2 is the same, except we apply Items 2 of Lemmas 33 and 26 respectively.

% Faithful official CCC2022 article16 Definition39, Theorem40 and Theorem43 with complete author proofs.
\section*{4.4 Case of PH (main promise-PH simulation)}
It is observed in [7, Corollary 17] that, for every $\ell\geq1$, $\Sigma_\ell^{\mathrm P}=\operatorname{lin}\Sigma_\ell^{\mathrm P}$, under the assumption that $\mathrm{Dist}\Sigma_{\ell+1}^{\mathrm P}\subseteq\mathrm{AvgP}$. We will show a similar result for the class $\operatorname{promise}\text{-}\Sigma_\ell^{\mathrm{BP}}$ defined below.
\paragraph{Definition 39 ($\operatorname{promise}\text{-}\Sigma_\ell^{\mathrm{BP}}$ and $\operatorname{promise}\text{-}\operatorname{lin}\Sigma_\ell^{\mathrm{BP}}$).} A promise problem $\Pi=(\Pi_Y,\Pi_N)\in\operatorname{promise}\text{-}\Sigma_\ell^{\mathrm{BP}}$ for $\ell\geq0$, if there is a $\Sigma_\ell\circ\mathrm{BP}\circ\mathrm P$ formula
\[\varphi(x)=\exists y_1\forall y_2\ldots Qy_\ell\ \mathrm{BP}_r\ R(x,y_1,y_2,\ldots,y_\ell,r),\]
for a polytime predicate $R$ and all string variables $y_1,\ldots,y_\ell,r$ of length polynomial in $|x|$, such that, for all $x\in\{0,1\}^n$,
\begin{align*}
x\in\Pi_Y&\Longrightarrow\exists y_1\forall y_2\ldots Qy_\ell\ \mathrm{BP}_{(2/3)}r\ R(x,y_1,y_2,\ldots,y_\ell,r),\text{ and}\\
x\in\Pi_N&\Longrightarrow\forall y_1\exists y_2\ldots\overline Qy_\ell\ \mathrm{BP}_{(2/3)}r\ \neg R(x,y_1,y_2,\ldots,y_\ell,r),
\end{align*}
where $\mathrm{BP}_{(\delta)}r\ R(-,r)$ means that $\Pr_r[R(-,r)]\geq\delta$. The constant $2/3$ in the definition above can be changed to $1-2^{-\operatorname{poly}(n)}$ by standard success amplification techniques (even if we start with any pair $\delta_Y>1-\delta_N+\varepsilon$ of thresholds, for some $\varepsilon\geq1/\operatorname{poly}(n)$, where $\delta_Y$ is the probability $R(-,r)$ accepts, and $\delta_N$ the probability $\neg R(-,r)$ accepts).
The promise class $\operatorname{promise}\text{-}\operatorname{lin}\Sigma_\ell^{\mathrm{BP}}$ is defined the same way as above, with an extra condition that, for some constant $a>0$, $|y_i|\leq an$ for all $1\leq i\leq\ell$ (but the randomness $r$ may still be of any $\operatorname{poly}(n)$ length).
\paragraph{Theorem 40 (Witness Compression for PH).} For any $\ell\geq0$, suppose $\mathrm{Dist}\Sigma_{\ell+2}^{\mathrm P}\subseteq\mathrm{AvgBPP}$. Then
\[\operatorname{promise}\text{-}\Sigma_\ell^{\mathrm{BP}}=\operatorname{promise}\text{-}\operatorname{lin}\Sigma_\ell^{\mathrm{BP}}.\]
\paragraph{Proof.} The proof is by induction on $\ell$. The base case of $\ell=0$ is trivially true.
For any $\ell>0$, suppose the claim is true for $\ell-1$. Consider an arbitrary promise problem $\Pi=(\Pi_Y,\Pi_N)\in\operatorname{promise}\text{-}\Sigma_\ell^{\mathrm{BP}}$. That is, for every $x\in\{0,1\}^n$,
\begin{align*}
x\in\Pi_Y&\Longrightarrow\exists y\ \Pi_{\ell-1}z\ \mathrm{BP}_{(1-2^{-n})}r\ R(x,y,z,r),\\
x\in\Pi_N&\Longrightarrow\forall y\ \Sigma_{\ell-1}z\ \mathrm{BP}_{(1-2^{-n})}r\ \neg R(x,y,z,r),
\end{align*}
for some polytime predicate $R$, where $z$ is a sequence of $\ell-1$ strings $z_1,\ldots,z_{\ell-1}$ under the corresponding $\ell-1$ quantifiers of the $\Pi_{\ell-1}$ pattern, and $|y|,|r|,|z_i|\leq q(n)$, for some polynomial $q$, for all $1\leq i\leq\ell-1$.
As in the proof of the inclusion $\mathrm{BPP}\subseteq\Sigma_2^{\mathrm P}\cap\Pi_2^{\mathrm P}$ [22], we replace the BP quantifier by the $\exists\forall$ pattern if $\ell$ is odd, or by the $\forall\exists$ pattern if $\ell$ is even. We get that there exists a polytime predicate $R'$ such that, for every $x\in\{0,1\}^n$,
\begin{align*}
x\in\Pi_Y&\Longrightarrow\exists y\ \Pi_{\ell-1}z\ Q_0z'\ Q_1z''\ R'(x,y,z,z',z''),\\
x\in\Pi_N&\Longrightarrow\forall y\ \Sigma_{\ell-1}z\ Q_1z'\ Q_0z''\ \neg R'(x,y,z,z',z''),
\end{align*}
where $z',z''$ have the lengths polynomial in $n$ and $Q_0Q_1=\exists\forall$ if $\ell$ is odd, and $Q_0Q_1=\forall\exists$ if $\ell$ is even.
Note that the above transformation shows that our promise problem $\Pi\in\operatorname{promise}\text{-}\Sigma_{\ell+1}^{\mathrm P}$. For any given $x\in\Pi_Y$, let $y_x$ be the lexicographically first string $y$ such that
\[\Pi_{\ell-1}z\ Q_0z'\ Q_1z''\ R'(x,y,z,z',z'').\tag{23}\]
\paragraph{Claim 41.} For $x$ and $y_x$ as above,
\[\Pi_{\ell-1}z\ \mathrm{BP}_{(1/\operatorname{poly}(n))}r\ R(x,y_x,z,r).\]
\paragraph{Proof.} Immediate from the properties of Lautemann's proof of $\mathrm{BPP}\subseteq\Sigma_2^{\mathrm P}$ [22]. Recall that in that proof, the subformula $\mathrm{BP}_r\ R(-,r)$ is replaced by
\[\exists u_1\ldots u_k\ \forall r\ \bigvee_{i=1}^kR(-,r\oplus u_i),\]
for some $k\leq\operatorname{poly}(|r|)$, where $|u_i|=|r|$, for all $1\leq i\leq k$, and $\oplus$ is the bit-wise XOR. If $\Pr_r[R(-,r)]<1/k$, then such a collection of $u_1,\ldots,u_k$ cannot exist. Indeed, each ``shifted'' predicate $R(-,r\oplus u_i)$, for $1\leq i\leq k$, accepts less than $1/k$ fraction of $r\in\{0,1\}^{|r|}$, and hence, by the union bound, the OR of any $k$ such shifts accepts less than $k\cdot(1/k)=1$ fraction of $r\in\{0,1\}^{|r|}$.\hfill$\triangleleft$
On the other hand, for every $x\in\Pi_N$ and every $y$, we have
\[\Sigma_{\ell-1}z\ \mathrm{BP}_{(1-2^{-n})}r\ \neg R(x,y,z,r).\]
Using Claim 41, by standard success amplification, for a random string $r'$ of length $q'(n)$, for some polynomial $q'$, and a polytime predicate $R^{\mathrm{boost}}$, we get for every $x\in\{0,1\}^n$ and $y_x$ as defined above,
\begin{align*}
x\in\Pi_Y&\Longrightarrow\Pi_{\ell-1}z\ \mathrm{BP}_{(1-2^{-n})}r'\ R^{\mathrm{boost}}(x,y_x,z,r'),\tag{24}\\
x\in\Pi_N&\Longrightarrow\forall y\ \Sigma_{\ell-1}z\ \mathrm{BP}_{(1-2^{-n})}r'\ \neg R^{\mathrm{boost}}(x,y,z,r').\tag{25}
\end{align*}
\paragraph{Claim 42.} For some universal constant $a>0$ and some polynomial $p'$, for every $n$-bit $x\in\Pi_Y$ with a witness $y_x$ as defined above, we have
\[\mathrm{rK}^{p'(n)}(y_x\mid x)\leq an.\]
\paragraph{Proof.} Using Equation (23), we get by Item 2 of Lemma 34 (since $\mathrm{Dist}\Sigma_3^{\mathrm P}\subseteq\mathrm{AvgBPP}$), that, for some polynomial $p$ and all large enough $t$,
\[\mathrm{rK}^{p(t)}(y_x\mid x)\leq\mathrm{rK}^t(x)-\mathrm{rK}^{p(t)}(x)+2\log p(t).\]
Since, for any $t$, $\mathrm{rK}^t(x)\leq O(n)$, the claim follows.\hfill$\triangleleft$
Consider the formula
\[\varphi(x)=\exists y'\ (|y'|\leq an)\ \Pi_{\ell-1}z\ \mathrm{BP}_w\ R''(x,y',z,w),\tag{26}\]
where the predicate $R''$ is computed by the following polytime algorithm $A$:
\begin{itemize}
\item On input $x,y',z$, and a string $w=w'r'$, where $w'$ is of length $p'(n)$, and $r'$ is of length $q'(n)$, first run the universal machine $U(y',x,w')$ for $p'(n)$ steps, getting an output $y$. Then run $R^{\mathrm{boost}}(x,y,z,r')$, accepting iff $R^{\mathrm{boost}}$ accepts.
\end{itemize}
We claim that Equation (26) is a correct $\Sigma_\ell^{\mathrm{BP}}$ formula for the promise problem $\Pi$. Indeed, for any $x\in\Pi_Y$, by Claim 42, there is a string $y'$ that is decompressed to $y_x$ with high probability over random strings $w'$ (say, with probability at least $3/4$). By Equation (24), $R^{\mathrm{boost}}(x,U(y',x,w'),z,r')$ accepts with probability at least $(3/4)(1-2^{-n})>2/3$ over strings $w=w'r'$. On the other hand, for any $x\in\Pi_N$, Equation (25) implies that
\[\forall y'\ \Sigma_{\ell-1}z\ \mathrm{BP}_{(1-2^{-n})}w\ \neg R''(x,y',z,w).\]
Next, using Equation (26), define a new promise problem $\Pi'=(\Pi'_Y,\Pi'_N)\in\operatorname{promise}\text{-}\Pi_{\ell-1}^{\mathrm{BP}}$:
\begin{align*}
\Pi'_Y&=\{(x,y')\mid|y'|\leq a|x|,\ \Pi_{\ell-1}z\ \mathrm{BP}_w\ R''(x,y',z,w)\},\\
\Pi'_N&=\{(x,y')\mid|y'|\leq a|x|,\ \Sigma_{\ell-1}z\ \mathrm{BP}_w\ \neg R''(x,y',z,w)\}.
\end{align*}
By the Inductive Hypothesis, $\Pi'\in\operatorname{promise}\text{-}\operatorname{lin}\Pi_{\ell-1}^{\mathrm{BP}}$, with all quantified binary strings of length $O(|x|+|y'|)\leq O(n)$. It follows that $\Pi\in\operatorname{promise}\text{-}\operatorname{lin}\Sigma_\ell^{\mathrm{BP}}$.\hfill$\triangleleft$

\paragraph{Theorem 43.} For any $\ell\geq0$,
\[\mathrm{Dist}\Sigma_{\ell+2}^{\mathrm P}\subseteq\mathrm{AvgBPP}\Longrightarrow\operatorname{promise}\text{-}\Sigma_\ell^{\mathrm P}\subseteq\operatorname{promise}\text{-}\mathrm{BPTIME}[2^{O(n/\log n)}].\]

\paragraph{Proof.} The proof is by induction on $\ell$. The base case of $\ell=0$ is trivially true. Consider any $\ell\geq1$, and assume the claim for $\ell-1$. Let $\Pi=(\Pi_Y,\Pi_N)\in\operatorname{promise}\text{-}\Sigma_\ell^{\mathrm P}$ be arbitrary. By Theorem 40, $\Pi\in\operatorname{promise}\text{-}\operatorname{lin}\Sigma_\ell^{\mathrm{BP}}$ via some formula
\[\varphi(x)=\exists y\ \Pi_{\ell-1}z\ \mathrm{BP}_r\ R(x,y,z,r),\]
with a polytime $R$, linearly-bounded $y$ and $z=(z_1,\ldots,z_{\ell-1})$ (under the $\Pi_{\ell-1}$ quantifiers), and a poly-bounded $r$.
Arguing as in the proof of Theorem 40 (see the proof of Claim 42), but using Lemma 31 (for the case of pK) to bound the computational depth of an input $x$ by $O(n/\log n)$ with a time bound $t\leq2^{O(n/\log n)}$, we get that, for some universal constant $a>0$, every $x\in\Pi_Y$ has a witness $y_x$ with
\[\mathrm{pK}^{2^{an/\log n}}(y_x\mid x)\leq an/\log n.\tag{27}\]
Applying success amplification to $R$, we get a new polytime predicate $R^{\mathrm{boost}}$ with poly-bounded randomness $r'$ such that we get for every $x\in\{0,1\}^n$ and $y_x$ as defined above,
\begin{align*}
x\in\Pi_Y&\Longrightarrow\Pi_{\ell-1}z\ \mathrm{BP}_{(1-2^{-n})}r'\ R^{\mathrm{boost}}(x,y_x,z,r'),\tag{28}\\
x\in\Pi_N&\Longrightarrow\forall y\ \Sigma_{\ell-1}z\ \mathrm{BP}_{(1-2^{-n})}r'\ \neg R^{\mathrm{boost}}(x,y,z,r').\tag{29}
\end{align*}
Define a promise problem $\Pi'=(\Pi'_Y,\Pi'_N)\in\operatorname{promise}\text{-}\Pi_{\ell-1}^{\mathrm{BP}}$:
\begin{align*}
\Pi'_Y&=\{(x,y)\mid\Pi_{\ell-1}z\ \mathrm{BP}_{r'}\ R^{\mathrm{boost}}(x,y,z,r')\},\\
\Pi'_N&=\{(x,y)\mid\Sigma_{\ell-1}z\ \mathrm{BP}_{r'}\ \neg R^{\mathrm{boost}}(x,y,z,r')\}.
\end{align*}
By the Inductive Hypothesis, $\Pi'\in\operatorname{promise}\text{-}\mathrm{BPTIME}[2^{O(n/\log n)}]$, where $|x|=n$, via some randomized algorithm $A$, with error probability at most $2^{-n}$.
We now solve $\Pi$ with the following randomized algorithm $B$:
\begin{itemize}
\item Given $x\in\{0,1\}^n$, choose a random string $w$ of length $2^{an/\log n}$, and then enumerate over all machines $M\in\{0,1\}^{an/\log n}$, and run $M(w,x)$ for $2^{an/\log n}$ steps, getting a candidate witness $y$ of length $O(n)$. Run $A(x,y)$, accepting iff $A$ accepts.
\end{itemize}
The time complexity of the described algorithm $B$ is clearly $2^{O(n/\log n)}$. For correctness, for every $x\in\Pi_Y$, with probability at least $3/4$ over $w$, our algorithm will consider a short description of a TM $M$ such that $M(w,x)=y_x$. Since $(x,y_x)\in\Pi'_Y$, we get that $A(x,y_x)$ will accept with probability at least $1-2^{-n}$ over its internal randomness. Hence, the algorithm $B$ accepts $x$ with probability at least $(3/4)(1-2^{-n})>2/3$. On the other hand, for every $x\in\Pi_N$, for every $w$ and all $M$, the string $y=M(w,x)$ is such that $(x,y)\in\Pi'_N$, and so $A(x,y)$ accepts with probability at most $2^{-n}$. It follows that the probability that $B(x)$ accepts in this case is
\[\mathbb E_{w,A}[\exists v\in\{0,1\}^{an/\log n}:A(x,U(v,x,w))]\leq2^{an/\log n}\cdot\mathbb E_w(2^{-n})\leq2^{-n/2},\]
for all sufficiently large $n$.\hfill$\triangleleft$

\begin{thebibliography}{45}
\bibitem[1]{ref1} Luis Filipe Coelho Antunes and Lance Fortnow. Worst-case running times for average-case algorithms. In Proceedings of the 24th Annual IEEE Conference on Computational Complexity, CCC 2009, Paris, France, 15-18 July 2009, pages 298–303. IEEE Computer Society, 2009. doi:10.1109/CCC.2009.12.
\bibitem[2]{ref2} Shai Ben-David, Benny Chor, Oded Goldreich, and Michael Luby. On the theory of average case complexity. J. Comput. Syst. Sci., 44(2):193–219, 1992. doi:10.1016/0022-0000(92)90019-F.
\bibitem[3]{ref3} Andrej Bogdanov and Luca Trevisan. Average-case complexity. Found. Trends Theor. Comput. Sci., 2(1), 2006. doi:10.1561/0400000004.
\bibitem[4]{ref4} Andrej Bogdanov and Luca Trevisan. On worst-case to average-case reductions for NP problems. SIAM J. Comput., 36(4):1119–1159, 2006. doi:10.1137/S0097539705446974.
\bibitem[5]{ref5} Harry Buhrman, Lance Fortnow, and Aduri Pavan. Some results on derandomization. Theory Comput. Syst., 38(2):211–227, 2005. doi:10.1007/s00224-004-1194-y.
\bibitem[6]{ref6} Lijie Chen, Shuichi Hirahara, and Neekon Vafa. Average-case hardness of NP and PH from worst-case fine-grained assumptions. In Mark Braverman, editor, 13th Innovations in Theoretical Computer Science Conference, ITCS 2022, January 31 - February 3, 2022, Berkeley, CA, USA, volume 215 of LIPIcs, pages 45:1–45:16. Schloss Dagstuhl - Leibniz-Zentrum für Informatik, 2022. doi:10.4230/LIPIcs.ITCS.2022.45.
\bibitem[7]{ref7} Halley Goldberg and Valentine Kabanets. A simpler proof of the worst-case to average-case reduction for polynomial hierarchy via symmetry of information. Electron. Colloquium Comput. Complex., 7:1–14, 2022. URL: https://eccc.weizmann.ac.il/report/2022/007/.
\bibitem[8]{ref8} Oded Goldreich. The Foundations of Cryptography - Volume 1: Basic Techniques. Cambridge University Press, 2001. doi:10.1017/CBO9780511546891.
\bibitem[9]{ref9} Oded Goldreich and Leonid A. Levin. A hard-core predicate for all one-way functions. In Proceedings of the Twenty-First Annual ACM Symposium on Theory of Computing, STOC ’89, pages 25–32, New York, NY, USA, 1989. Association for Computing Machinery. doi:10.1145/73007.73010.
\bibitem[10]{ref10} Shuichi Hirahara. Characterizing average-case complexity of PH by worst-case meta-complexity. In Sandy Irani, editor, 61st IEEE Annual Symposium on Foundations of Computer Science, FOCS 2020, Durham, NC, USA, November 16-19, 2020, pages 50–60. IEEE, 2020. doi:10.1109/FOCS46700.2020.00014.
\bibitem[11]{ref11} Shuichi Hirahara. Average-case hardness of NP from exponential worst-case hardness assump- tions. In Samir Khuller and Virginia Vassilevska Williams, editors, STOC ’21: 53rd Annual ACM SIGACT Symposium on Theory of Computing, Virtual Event, Italy, June 21-25, 2021, pages 292–302. ACM, 2021. doi:10.1145/3406325.3451065.
\bibitem[12]{ref12} Shuichi Hirahara. Symmetry of information in heuristica. Manuscript, 2021.
\bibitem[13]{ref13} Shuichi Hirahara and Mikito Nanashima. On worst-case learning in relativized heuristica. In Symposium on Foundations of Computer Science (FOCS), 2021.
\bibitem[14]{ref14} Shuichi Hirahara and Rahul Santhanam. Errorless versus error-prone average-case complexity. In Mark Braverman, editor, 13th Innovations in Theoretical Computer Science Conference, ITCS 2022, January 31 - February 3, 2022, Berkeley, CA, USA, volume 215 of LIPIcs, pages 84:1–84:23. Schloss Dagstuhl - Leibniz-Zentrum für Informatik, 2022. doi:10.4230/LIPIcs.ITCS.2022.84.
\bibitem[15]{ref15} Shuichi Hirahara and Rahul Santhanam. Excluding PH pessiland. In Mark Braverman, editor, 13th Innovations in Theoretical Computer Science Conference, ITCS 2022, January 31 - February 3, 2022, Berkeley, CA, USA, volume 215 of LIPIcs, pages 85:1–85:25. Schloss Dagstuhl - Leibniz-Zentrum für Informatik, 2022. doi:10.4230/LIPIcs.ITCS.2022.85.
\bibitem[16]{ref16} Russell Impagliazzo. A personal view of average-case complexity. In Proceedings of the Tenth Annual Structure in Complexity Theory Conference, Minneapolis, Minnesota, USA, June 19-22, 1995, pages 134–147, 1995. doi:10.1109/SCT.1995.514853.
\bibitem[17]{ref17} Russell Impagliazzo. Relativized separations of worst-case and average-case complexities for NP. In Proceedings of the 26th Annual IEEE Conference on Computational Complexity, CCC 2011, San Jose, California, USA, June 8-10, 2011, pages 104–114, 2011. doi:10.1109/CCC.2011.34.
\bibitem[18]{ref18} Russell Impagliazzo and Leonid A. Levin. No better ways to generate hard NP instances than picking uniformly at random. In 31st Annual Symposium on Foundations of Computer Science, St. Louis, Missouri, USA, October 22-24, 1990, Volume II, pages 812–821. IEEE Computer Society, 1990. doi:10.1109/FSCS.1990.89604.
\bibitem[19]{ref19} Russell Impagliazzo and Avi Wigderson. P = BPP if E requires exponential circuits: De- randomizing the XOR lemma. In Frank Thomson Leighton and Peter W. Shor, editors, Proceedings of the Twenty-Ninth Annual ACM Symposium on the Theory of Computing, El Paso, Texas, USA, May 4-6, 1997, pages 220–229. ACM, 1997. doi:10.1145/258533.258590.
\bibitem[20]{ref20} Michael J. Kearns, Robert E. Schapire, and Linda Sellie. Toward efficient agnostic learning. Mach. Learn., 17(2-3):115–141, 1994. doi:10.1007/BF00993468.
\bibitem[21]{ref21} Pravesh K. Kothari and Roi Livni. Improper learning by refuting. In Anna R. Karlin, editor, 9th Innovations in Theoretical Computer Science Conference, ITCS 2018, January 11-14, 2018, Cambridge, MA, USA, volume 94 of LIPIcs, pages 55:1–55:10. Schloss Dagstuhl - Leibniz-Zentrum für Informatik, 2018. doi:10.4230/LIPIcs.ITCS.2018.55.
\bibitem[22]{ref22} Clemens Lautemann. BPP and the polynomial hierarchy. Inf. Process. Lett., 17(4):215–217, 1983. doi:10.1016/0020-0190(83)90044-3.
\bibitem[23]{ref23} Troy Lee. Kolmogorov complexity and formula lower bounds. PhD thesis, University of Amsterdam, 2006.
\bibitem[24]{ref24} Leonid A. Levin. Laws of information conservation (nongrowth) and aspects of the foundation of probability theory. Problemy Peredachi Informatsii, 10(3):30–35, 1974.
\bibitem[25]{ref25} Leonid A. Levin. Average case complete problems. SIAM J. Comput., 15(1):285–286, 1986. doi:10.1137/0215020.
\bibitem[26]{ref26} Ming Li and Paul M. B. Vitányi. An introduction to Kolmogorov complexity and its applications. Springer, 2019.
\bibitem[27]{ref27} Ming Li and Paul M. B. Vitányi. An Introduction to Kolmogorov Complexity and Its Applications, 4th Edition. Texts in Computer Science. Springer, 2019. doi:10.1007/978-3-030-11298-1.
\bibitem[28]{ref28} Yanyi Liu and Rafael Pass. On one-way functions and kolmogorov complexity. In 61st IEEE Annual Symposium on Foundations of Computer Science, FOCS 2020, Durham, NC, USA, November 16-19, 2020, pages 1243–1254, 2020. doi:10.1109/FOCS46700.2020.00118.
\bibitem[29]{ref29} Yanyi Liu and Rafael Pass. On one-way functions from np-complete problems. Electron. Colloquium Comput. Complex., page 59, 2021. URL: https://eccc.weizmann.ac.il/report/2021/059.
\bibitem[30]{ref30} Yanyi Liu and Rafael Pass. On the possibility of basing cryptography on $\mathrm{EXP}\neq\mathrm{BPP}$. In Advances in Cryptology - CRYPTO 2021 - 41st Annual International Cryptology Conference, CRYPTO 2021, Virtual Event, August 16-20, 2021, Proceedings, Part I, pages 11–40, 2021. doi:10.1007/978-3-030-84242-0\_2.
\bibitem[31]{ref31} Luc Longpré and Osamu Watanabe. On symmetry of information and polynomial time invertibility. Inf. Comput., 121(1):14–22, 1995. doi:10.1006/inco.1995.1120.
\bibitem[32]{ref32} Zhenjian Lu and Igor C. Oliveira. An efficient coding theorem via probabilistic representations and its applications. In Nikhil Bansal, Emanuela Merelli, and James Worrell, editors, 48th International Colloquium on Automata, Languages, and Programming, ICALP 2021, July 12-16, 2021, Glasgow, Scotland (Virtual Conference), volume 198 of LIPIcs, pages 94:1–94:20. Schloss Dagstuhl - Leibniz-Zentrum für Informatik, 2021. doi:10.4230/LIPIcs.ICALP.2021.94.
\bibitem[33]{ref33} Zhenjian Lu, Igor C. Oliveira, and Rahul Santhanam. Pseudodeterministic algorithms and the structure of probabilistic time. In Samir Khuller and Virginia Vassilevska Williams, editors, STOC ’21: 53rd Annual ACM SIGACT Symposium on Theory of Computing, Virtual Event, Italy, June 21-25, 2021, pages 303–316. ACM, 2021. doi:10.1145/3406325.3451085.
\bibitem[34]{ref34} Zhenjian Lu, Igor C. Oliveira, and Marius Zimand. Optimal coding theorems in time-bounded kolmogorov complexity. In 49th International Colloquium on Automata, Languages and Programming, 2022.
\bibitem[35]{ref35} Peter Bro Miltersen, N. V. Vinodchandran, and Osamu Watanabe. Super-polynomial versus half-exponential circuit size in the exponential hierarchy. In Takao Asano, Hiroshi Imai, D. T. Lee, Shin-Ichi Nakano, and Takeshi Tokuyama, editors, Computing and Combinatorics, 5th Annual International Conference, COCOON ’99, Tokyo, Japan, July 26-28, 1999, Proceedings, volume 1627 of Lecture Notes in Computer Science, pages 210–220. Springer, 1999. doi:10.1007/3-540-48686-0\_21.
\bibitem[36]{ref36} Ashish V. Naik, Kenneth W. Regan, and D. Sivakumar. On quasilinear-time complexity theory. Theor. Comput. Sci., 148(2):325–349, 1995. doi:10.1016/0304-3975(95)00031-Q.
\bibitem[37]{ref37} Noam Nisan and Avi Wigderson. Hardness vs randomness. J. Comput. Syst. Sci., 49(2):149–167, 1994. doi:10.1016/S0022-0000(05)80043-1.
\bibitem[38]{ref38} Igor C. Oliveira. Randomness and intractability in kolmogorov complexity. In Christel Baier, Ioannis Chatzigiannakis, Paola Flocchini, and Stefano Leonardi, editors, 46th International Colloquium on Automata, Languages, and Programming, ICALP 2019, July 9-12, 2019, Patras, Greece, volume 132 of LIPIcs, pages 32:1–32:14. Schloss Dagstuhl - Leibniz-Zentrum für Informatik, 2019. doi:10.4230/LIPIcs.ICALP.2019.32.
\bibitem[39]{ref39} Igor C. Oliveira and Rahul Santhanam. Pseudodeterministic constructions in subexponential time. In Hamed Hatami, Pierre McKenzie, and Valerie King, editors, Proceedings of the 49th Annual ACM SIGACT Symposium on Theory of Computing, STOC 2017, Montreal, QC, Canada, June 19-23, 2017, pages 665–677. ACM, 2017. doi:10.1145/3055399.3055500.
\bibitem[40]{ref40} Hanlin Ren and Rahul Santhanam. Hardness of KT characterizes parallel cryptography. In 36th Computational Complexity Conference, CCC 2021, July 20-23, 2021, Toronto, Ontario, Canada (Virtual Conference), pages 35:1–35:58, 2021. doi:10.4230/LIPIcs.CCC.2021.35.
\bibitem[41]{ref41} Shai Shalev-Shwartz and Shai Ben-David. Understanding Machine Learning - From Theory to Algorithms. Cambridge University Press, 2014. URL: http://www.cambridge.org/de/academic/subjects/computer-science/pattern-recognition-and-machine-learning/ understanding-machine-learning-theory-algorithms.
\bibitem[42]{ref42} Ronen Shaltiel and Christopher Umans. Simple extractors for all min-entropies and a new pseudorandom generator. J. ACM, 52(2):172–216, 2005. doi:10.1145/1059513.1059516.
\bibitem[43]{ref43} Luca Trevisan and Salil P. Vadhan. Pseudorandomness and average-case complexity via uniform reductions. Computational Complexity, 16(4):331–364, 2007.
\bibitem[44]{ref44} Salil P. Vadhan. On learning vs. refutation. In Satyen Kale and Ohad Shamir, editors, Proceed- ings of the 30th Conference on Learning Theory, COLT 2017, Amsterdam, The Netherlands, 7-10 July 2017, volume 65 of Proceedings of Machine Learning Research, pages 1835–1848. PMLR, 2017. URL: http://proceedings.mlr.press/v65/vadhan17a.html.
\bibitem[45]{ref45} Emanuele Viola. On constructing parallel pseudorandom generators from one-way functions. In 20th Annual IEEE Conference on Computational Complexity (CCC 2005), 11-15 June 2005, San Jose, CA, USA, pages 183–197, 2005. doi:10.1109/CCC.2005.16.
\end{thebibliography}
\end{document}
